\documentclass{article}
\usepackage[UTF8]{ctex}
\usepackage{amsmath,amssymb,amsfonts}
\usepackage{enumitem}
\usepackage{geometry}
\geometry{a4paper, margin=1in}
% === 错题本排版优化（不参与题目标识） ===
\setlength{\parindent}{0pt}
\setlength{\parskip}{0.35em plus 0.10em minus 0.08em}
\setlength{\emergencystretch}{3em}
\setlist{itemsep=0.25em, topsep=0.40em, parsep=0pt}
\allowdisplaybreaks[2]
% === 错题本排版优化结束 ===
\begin{document}
\section*{华东理工大学2024年数学分析真题}
\section*{53．华东理工大学2024年数学分析试卷（回忆版）}

\subsection{题目1}
一．求极限．

\begin{itemize}
  \item[（1）] \(\lim _{n \rightarrow \infty} \sqrt[n]{2 \sin ^{2} n+\cos ^{2} n}\) ．
  \item[（2）] \(\lim _{x \rightarrow 0} \frac{\cos x-e^{-\frac{x^{2}}{2}}}{x^{4}}\) ．
\end{itemize}

\subsection{题目2}
二．求重积分 \(\iint_{D} e^{\max \left\{x^{2}, y^{2}\right\}} \mathrm{d} x \mathrm{~d} y\) ，其中 \(D:[0,1] \times[0,1]\) ．

\subsection{题目3}
三．设 \(f(x)\) 是 \([1,+\infty)\) 上单调递减的连续函数．证明：

\begin{itemize}
  \item[（1）] \(\int_{1}^{+\infty} f(x) \mathrm{d} x\) 收敛时，有 \(\lim _{x \rightarrow+\infty} x f(x)=0\) 。
  \item[（2）] 举例说明：存在不单调递减的函数 \(f(x)\) 使得 \(\int_{1}^{+\infty} f(x) \mathrm{d} x\) 收敛，但 \(f(x)\) 在 \(+\infty\) 的极限不存在．
\end{itemize}

\subsection{题目4}
四．求 \(f(x, y, z)=\ln x+2 \ln y+3 \ln z\) 在 \(x^{2}+y^{2}+z^{2}=r^{2}(x, y, z>0, r>0)\) 下的最大值，并证明对任的正数 \(a, b, c\) ，都有

\[
a b^{2} a^{3} \leq 108\left(\frac{a+b+c}{6}\right)^{6}
\]

\subsection{题目5}
五．已知 \(a_{n}>0\) ，证明：

\[
\lim _{n \rightarrow \infty} \frac{a_{n+1}}{a_{n}} \leq \lim _{n \rightarrow \infty} \sqrt[n]{a_{n}} \leq \lim _{n \rightarrow \infty} \sqrt[n]{a_{n}} \leq \lim _{n \rightarrow \infty} \frac{a_{n+1}}{a_{n}}
\]

\subsection{题目6}
六．设函数列 \(\left\{S_{n}(x)\right\}\) 在 \([0, a]\) 上连续，且 \(S_{n}(x)=\int_{0}^{x} S_{n-1}(t) \mathrm{d} t\) ，证明：\(\left\{S_{n}(x)\right\}\) 在 \([0, a]\) 上一致收敛。

\subsection{题目7}
七．求曲面积分

\[
I=\iint_{\sum} \frac{x \mathrm{~d} y \mathrm{~d} z+y \mathrm{~d} z \mathrm{~d} x+z \mathrm{~d} x \mathrm{~d} y}{\left(x^{2}+y^{2}+z^{2}\right)^{\frac{3}{2}}} .
\]

其中 \(\sum\) 是立方体 \(|x| \leq 1,|y| \leq 2,|z| \leq 3\) 的表面，取外侧。

\subsection{题目8}
八．证明函数

\[
f(x)= \begin{cases}x^{\alpha} \cos \frac{1}{x}, & x \in\left(0, \frac{\pi}{2}\right] ; \\ \left(x-\frac{\pi}{2}\right)^{2024} e^{-x}, & x \in\left(\frac{\pi}{2},+\infty\right)\end{cases}
\]

在 \((0,+\infty)\) 上一致连续，其中 \(\alpha>0\) ．

\subsection{题目9}
九．讨论 \(\int_{1}^{+\infty} \cos \left(x^{p}\right) \mathrm{d} x\) 的绝对收敛性和条件收敛性．
\end{document}
